\section{Theory}

\begin{theorem}[Node quantile split]
\label{thm:quantile}
Let $X\sim P_{\calN}$ be the feature distribution at node $\calN$.  Let
$q\in\Sphere^{d-1}$ and suppose $T=q^\top \tilde X$ has continuous CDF $F$.
For $\rho\in(0,1)$, define $\tau=F^{-1}(\rho)$.  Then
\begin{equation}
  \Prob(T\le \tau)=\rho,\qquad \Prob(T>\tau)=1-\rho .
\end{equation}
\end{theorem}

\begin{proof}
This is the defining property of the $\rho$-quantile for a continuous scalar
random variable.
\end{proof}

\begin{theorem}[Recursive uniform dispatch]
\label{thm:recursive}
Suppose every internal node $\calN$ in the SPR tree splits its feature mass in
proportion to the number of experts assigned to its two children:
\begin{equation}
  \Prob(C_L\mid C_{\calN})=\frac{|\calE_L|}{|\calE_{\calN}|},
  \qquad
  \Prob(C_R\mid C_{\calN})=\frac{|\calE_R|}{|\calE_{\calN}|}.
\end{equation}
If the root contains all $E$ experts and each leaf contains one expert, then
each leaf cell has mass $1/E$.
\end{theorem}

\begin{proof}
For any node $\calN$, prove by induction that
$\Prob(C_{\calN})=|\calE_{\calN}|/E$.  The root satisfies this because its
mass is one and it contains $E$ experts.  If the claim holds for a parent node,
then
\begin{equation}
  \Prob(C_L)
  =\Prob(C_{\calN})\Prob(C_L\mid C_{\calN})
  =\frac{|\calE_{\calN}|}{E}
   \frac{|\calE_L|}{|\calE_{\calN}|}
  =\frac{|\calE_L|}{E},
\end{equation}
and the right child is identical.  At a leaf, $|\calE_e|=1$, hence
$\Prob(C_e)=1/E$.
\end{proof}

\begin{proposition}[SPR is not batch capacity assignment]
After calibration, SPR evaluates stored directions and thresholds along one
root-to-leaf path.  The route of $x$ does not depend on the other examples in
the current batch or on current expert counts.
\end{proposition}

\paragraph{Finite calibration.}
Empirical quantiles give exact or near-exact balance on the calibration set.
Held-out balance depends on calibration size and distribution match.  For one
scalar split, empirical CDF concentration gives the standard
$O(\sqrt{\log(1/\delta)/n_{\calN}})$ scale at node sample size $n_{\calN}$
\cite{dvoretzky1956asymptotic}.  Lower tree nodes have fewer samples, so
calibration size is a real experimental variable.

\paragraph{Training margins.}
Uniform dispatch at calibration time is only the initialization problem.  For a
tree router, we define the margin of $x$ as the closest distance to a threshold
along the selected path:
\begin{equation}
  m_{\mathrm{tree}}(x)=
  \min_{\calN\in\mathrm{path}(x)}
  |q_{\calN}^{\top}\tilde x-\tau_{\calN}|.
\end{equation}
Small lower-tail margins predict route flips under representation drift, so the
training experiments report loss, usage, route-feature NMI, flip rate, and
margin lower tails.
